\documentclass[10pt,a4paper]{amsart}
\usepackage[margin=19mm]{geometry}
\usepackage{amsmath,amssymb,mathtools}
\usepackage[scr=rsfs,cal=euler]{mathalfa}
\usepackage{xfrac}
\usepackage[unicode,hidelinks,bookmarks=false]{hyperref}
\newcommand{\ie}{i.e.\ }
\newcommand{\cf}{cf.\ }
\newcommand{\resp}{respectively\ }
\newcommand{\Subsection}{\subsection}
\providecommand{\nobreakdash}{\nobreak\hbox{-}}
\setlength{\emergencystretch}{2em}
\multlinegap0pt
\usepackage{amssymb}
\usepackage{enumerate}
\usepackage[normalem]{ulem}
\usepackage{xcolor}
\newtheorem{proposition}{Proposition}
\title{Some aspects of Equivariant homotopy dense subset}
\author{Sergey A. Antonyan, Luis A. Mart\'inez-S\'anchez}
\date{Source: arXiv:2605.24141v1 (CC BY 4.0)}
\begin{document}
\maketitle

\noindent\emph{Source and licence.} Some aspects of Equivariant homotopy dense subset. Version arXiv:2605.24141v1 (CC BY 4.0), distributed by arXiv under the Creative Commons Attribution 4.0 International licence, http://creativecommons.org/licenses/by/4.0/, as recorded in the arXiv licence metadata for this exact version and shown on the abstract page https://arxiv.org/abs/2605.24141v1. Full licence notice: \texttt{licenses/arxiv-2605.24141-CC-BY-4.0.txt}.\par\smallskip

\noindent\textit{Editorial scope.} Counted unit: the article's proposition Cone(X)semilattice (source0.tex lines 632--643). The article's complete original proof of it is delivered with it (source0.tex lines 645--683, one \texttt{proof} environment of 39 lines). No mathematics is written, repaired or re-derived: the statement and the proof are the article's own lines, in the article's own notation, and no retained line is substituted or re-flowed.  No further source-local context is needed: every cross-reference of every retained line resolves inside the two delivered spans.  Nothing was omitted from any retained span, so the package carries no omission record.  No retained line contains a citation, so the wrapper carries no bibliography and no bibliography entry.  No retained line contains a figure environment, an image-inclusion command or drawing code, so no image is delivered and the package declares no asset dependency.  Editorial formatting only, in addition: an AMS \texttt{amsart} preamble that restates the article's own declarations and macros used by the retained lines, namely the environments \texttt{proposition} on separate counters, together with the standard packages those lines need.  The retained environments print in this document as Proposition 1 in order of appearance, whereas the article prints the same items with the numbering of its own full document.  The complete native source of the article as distributed in the arXiv e-print for this exact version is bundled unchanged as \texttt{sources/201252/source0.tex}.
	\begin{proposition}\label{Cone(X)semilattice}
		\rm{Let $S$ be a $G$-semilattice, and let $\mathrm{Cone}(S)$ be the cone of $S$ equipped with the $G$-action $g\cdot sx=s(gx)$. Consider the map 
		$$\cdot: \mathrm{Cone}(S)\times \mathrm{Cone}(S)\rightarrow \mathrm{Cone}(S)$$
		 defined by 
		 $$sx\cdot ty=\mbox{min}\{s,t\}(xy).$$ 
		 Then the following conditions are satisfied:
			
			\begin{enumerate}
				\item [(i)] $(\mathrm{Cone}(S),\cdot)$ is a $G$-semilattice.
				\item [(ii)] If $S$ is a Lawson $G$-semilattice, then so is $\mathrm{Cone}(S)$.
		\end{enumerate}}
	\end{proposition}

	\begin{proof}
		
		(i) We first prove that the multiplication is continuous. To this end, let $\eta: I\times S\rightarrow \mathrm{Cone}(S)$ denote the quotient map. Take $sx, ty\in \mathrm{Cone}(S)$, and let $W$ be an open neighborhood of $sx\cdot ty=\text{min}\{s,t\}(xy)$ in $\mathrm{Cone}(S)$. We consider two cases:
		\begin{itemize}
			\item \textbf{Case 1}: $sx\neq \theta\neq ty$. In this case, we have $\text{min}\{s,t\}(xy)\neq \theta$, and thus, $\eta^{-1}(W)$ is an open neighborhood of $(\text{min}\{s,t\}, xy)$ in $I\times S$. Therefore, there exist $\delta>0$ and neighborhoods $V_x$, $V_y$ of $x$ and $y$, respectively, such that $0<s-\delta<s+\delta<1$, $0<t-\delta<t+\delta<1$ and $\text{min}\big((s-\delta,s+\delta)\times (t-\delta,t+\delta)\big)\times V_xV_y\subset \eta^{-1}(W)$. 
			
			Now, set
			
			\begin{center}
				$U_x=\eta\big((s-\delta,s+\delta)\times V_x\big)\quad$ and $\quad U_y=\eta\big((t-\delta,t+\delta)\times V_y\big)$.
			\end{center}
			
			Then $U_x$ and $U_y$ are neighborhoods of $sx$ and $ty$ in $\mathrm{Cone}(S)$, respectively, and satisfy $U_x\cdot U_y\subset W$. Hence, the multiplication is continuous at $(sx,ty)$.
			
			\item \textbf{Case 2}: $sx=\theta$ or $ty=\theta$. Without loss of generality, suppose that $sx=\theta$. By definition of the multiplication, we have $sx\cdot ty=\theta$. Hence, there exists $\varepsilon > 0$ such that $[0,\varepsilon)\times S\subset \eta^{-1}(W)$. Notice that $\eta([0,\varepsilon)\times S)$ and $\mathrm{Cone}(S)$ are open neighborhoods of $sx$ and $ty$, respectively, and moreover $\eta([0,\varepsilon)\times S)\cdot \mathrm{Cone}(S)\subset W$. Thus, the multiplication is continuous at $(sx,ty)$.
		\end{itemize}
		
		
		
		On the other hand, the multiplication is associative. Indeed, take $rx,sy,tz\in \mathrm{Cone}(S)$, and let us prove that $rx\cdot(sy\cdot tz)=(rx\cdot sy)\cdot tz$. If $\theta \in \{rx, sy, tz\}$, then both sides equal $\theta$. Otherwise,
		\begin{center}
			$rx\cdot(sy\cdot tz)=\mbox{min}\{r,\mbox{min}\{s,t\}\}(x(yz))=\mbox{min}\{\mbox{min}\{r,s\},t\}((xy)z)=(rx\cdot sy)\cdot tz$.  
		\end{center}
		
		Finally, we verify equivariance. Given $g\in G$ and $rx,sy\in \mathrm{Cone}(S)$, we have
		
		\begin{center}
			$g\cdot (rx\cdot sy)=g\cdot (\mbox{min}\{r,s\}(xy))=\mbox{min}\{r,s\}(g\cdot (xy))=\mbox{min}\{r,s\}((gx)\cdot (gy))=(g\cdot rx)(g\cdot sy)$.
		\end{center}
		Therefore, $(\mathrm{Cone}(S),\cdot)$ is a topological $G$-semilattice.   
		
		(ii) It suffices to prove that any point of $\mathrm{Cone}(S)$ admits a neighborhood basis consisting of subsemilattices of $\mathrm{Cone}(S)$. Let $sx\in \mathrm{Cone}(S)$ and let $V$ be an open neighborhood of $sx$.
		
		\begin{itemize}
			\item \textbf{Case 1}: $sx\neq \theta$. In this case, $s > 0$, and hence $\eta^{-1}(V)$ is an open neighborhood of $(s,x)$ in $I\times S$. Since $S$ is a Lawson semilattice, there exists a neighborhood $U$ of $x$ in $S$ such that $U$ is a semilattice and $(s-\varepsilon,s+\varepsilon)\times U\subset \eta^{-1}(V)$, for some $\varepsilon>0$ with  $0<s-\varepsilon<s+\varepsilon<1$. Then $\eta\big((s-\varepsilon,s+\varepsilon)\times U\big)$ is a neighborhood of $sx$, contained in $V$, which is itself a subsemilattice of $\mathrm{Cone}(S)$.
			
			\item \textbf{Case 2}: $sx=\theta$. In this case, there exists $\delta>0$ such that $[0,\delta)\times S\subset \eta^{-1}(V)$. Thus, $U=\eta([0,\delta)\times S)$ is the desired neighborhood of $sx$.
		\end{itemize}
	\end{proof}

\end{document}
